% =============================================================================
% Standalone note: which viscous projectors are Korn-compatible?
%
% Created 2026-07-29.  It carries the material that was removed from the article's
% Sec. 2.1 ("The viscous projector") when that subsection was slimmed down: the
% SHARP characterization of the Korn compatibility (V2), the ellipticity constant
% c_Pi, the Plancherel proof of sufficiency, the oscillating-profile proof of
% necessity, and the projectors that are admissible without lying in the family
% covered by the article's sufficient criterion.
%
% The article needs only the sufficient criterion, because all four operators of
% interest satisfy it.  This note supplies the converse half and the boundary of
% the class, and is the reference for the claim -- made in the article's App. B,
% rem:ftGenericPi -- that for a general admissible projector the discarded O(1)
% viscous factor lies in [2 c_Pi^2, 2] rather than in the uniform window [1,2].
%
% All the algebra below is machine-checked by
% proof_verification/sympy/projector_algebra_verification.py (sections C and E)
% and, for the abstract chaining, by proof_verification/coq-formal/ViscousProjector.v.
%
% LaTeX build hygiene: latexmkrc routes every intermediate into
% "latex compilation/<basename>/"; only the PDF is left beside this source.
% =============================================================================
\documentclass[11pt]{article}
\usepackage[margin=2.7cm]{geometry}
\usepackage{amsmath,amssymb,amsthm,mathtools}
\usepackage{enumitem}
\usepackage[colorlinks=true,allcolors=blue]{hyperref}
\usepackage[capitalize,noabbrev]{cleveref}

% (V1)/(V2) are enumitem items with a custom ref format.  Without these four lines
% cleveref prefixes the enumi type name and prints the clunky "Item (V1)".  Same
% suppression as the article's preamble, so the two documents read alike.
\crefformat{enumi}{#2#1#3}
\crefrangeformat{enumi}{#3#1#4--#5#2#6}
\Crefformat{enumi}{Property~#2#1#3}
\crefmultiformat{enumi}{#2#1#3}{ and~#2#1#3}{, #2#1#3}{, and~#2#1#3}
\Crefmultiformat{enumi}{Properties~#2#1#3}{ and~#2#1#3}{, #2#1#3}{, and~#2#1#3}

% ------------------------- theorem environments -------------------------
\theoremstyle{plain}
\newtheorem{theorem}{Theorem}[section]
\newtheorem{lemma}[theorem]{Lemma}
\newtheorem{proposition}[theorem]{Proposition}
\newtheorem{corollary}[theorem]{Corollary}
\theoremstyle{definition}
\newtheorem{definition}[theorem]{Definition}
\theoremstyle{remark}
\newtheorem{remark}[theorem]{Remark}
\newtheorem{example}[theorem]{Example}

% ------------------------- notation -------------------------
% Matches the article: \ViscProj is the GENERIC viscous projector, \DSPi its
% deviatoric-symmetric instance, \SPi and \DPi the two elementary extractors.
\newcommand{\ViscProj}{\overset{\scriptscriptstyle{\mathrm{v}}}{\Pi}}
\newcommand{\DSPi}{\overset{\scriptscriptstyle{\mathrm{DS}}}{\Pi}}
\newcommand{\DPi}{\overset{\scriptscriptstyle{\mathrm{D}}}{\Pi}}
\newcommand{\SPi}{\overset{\scriptscriptstyle{\mathrm{S}}}{\Pi}}
\newcommand{\bv}{\boldsymbol{v}}
\newcommand{\bw}{\boldsymbol{w}}
\newcommand{\bk}{\boldsymbol{k}}
\newcommand{\bx}{\boldsymbol{x}}
\newcommand{\be}{\boldsymbol{e}}
\newcommand{\dd}{\,\mathrm{d}}
\newcommand{\CK}{C_{\mathrm{K}}}
\newcommand{\cP}{c_{\Pi}}
\DeclarePairedDelimiter{\norm}{\lVert}{\rVert}

\title{Korn{-}compatible viscous projectors:\\ a sharp characterization}
\author{Companion note to~\cite{paper}}
\date{}

\begin{document}
\maketitle

\begin{abstract}
The stabilized formulation of the accompanying article~\cite{paper} is written for a
\emph{viscous projector} $\ViscProj$ subject to two requirements: that it be a pointwise,
constant{-}coefficient orthogonal projection of second{-}order tensors, and that it be
\emph{Korn{-}compatible}, i.e.\ that $\norm{\nabla\bv}\le\CK\norm{\ViscProj\nabla\bv}$ on
$H^1_0(\Omega)^d$. The article verifies the second requirement through a sufficient criterion
--- containment of the deviatoric{-}symmetric tensors in the range --- which covers the four
operators it uses. This note supplies what the article does not need: the criterion is
\emph{not} necessary, and Korn compatibility admits a sharp, purely algebraic and
domain{-}independent characterization, namely that the kernel of the projector contain no
nonzero rank{-}one tensor. We prove both implications, compute the resulting optimal constant,
exhibit an admissible projector lying outside the article's family, identify the two
elementary projectors that are \emph{not} admissible, and draw the one consequence the
article records: for a general admissible projector the $O(1)$ viscous factor discarded in the
design of the stabilization parameters is no longer confined to the uniform window $[1,2]$.
\end{abstract}

\section{Setting and notation}\label{sec:setting}

Let $\Omega\subset\mathbb{R}^d$, $d\ge2$, be a bounded domain. Second{-}order tensors carry the
Frobenius inner product $\mathbf{A}\!:\!\mathbf{B}=A_{ij}B_{ij}$ and the induced norm
$|\mathbf{A}|$; $\norm{\cdot}$ denotes the $L^2(\Omega)$ norm of the corresponding field.
Following~\cite{paper}, a linear map $\Pi$ on second{-}order tensors, acting pointwise and with
constant coefficients, is called \emph{admissible} when it satisfies

\begin{enumerate}[label=\textup{(V\arabic*)},ref=\textup{(V\arabic*)},leftmargin=3.4em]
  \item\label{V:orth} \emph{Orthogonality.} $\Pi$ is an orthogonal projection for the Frobenius
  inner product: $\Pi\Pi=\Pi$ and $(\Pi\mathbf{A})\!:\!\mathbf{B}=\mathbf{A}\!:\!(\Pi\mathbf{B})$
  for all $\mathbf{A},\mathbf{B}$. In particular $|\Pi\mathbf{T}|\le|\mathbf{T}|$ pointwise.
  \item\label{V:korn} \emph{Korn compatibility.} There is $\CK\ge1$, independent of the mesh and
  of the physical parameters, with
  \begin{equation}\label{eq:PiKorn}
    \norm{\nabla\bv}\ \le\ \CK\,\norm{\Pi\nabla\bv}
    \qquad\forall\,\bv\in H^1_0(\Omega)^d .
  \end{equation}
\end{enumerate}

\noindent Two elementary extractors generate the operators used in~\cite{paper}:
\begin{equation}\label{eq:extractors}
  \SPi\mathbf{T}\coloneqq\operatorname{sym}\mathbf{T}
  =\tfrac12\bigl(\mathbf{T}+\mathbf{T}^{\mathsf T}\bigr),
  \qquad
  \DPi\mathbf{T}\coloneqq\mathbf{T}-\tfrac1d(\operatorname{tr}\mathbf{T})\,\mathbf{I},
\end{equation}
which commute, so that $\DSPi\coloneqq\DPi\SPi=\SPi\DPi$ extracts the deviatoric{-}symmetric
part. The article's family is
\begin{equation}\label{eq:family}
  \mathcal{F}\coloneqq\bigl\{\,\mathbb{I},\ \SPi,\ \DPi,\ \DSPi\,\bigr\},
\end{equation}
and its sufficient criterion reads as follows.

\begin{proposition}[Sufficient criterion; Lemma 2.1 of~\cite{paper}]\label{prop:criterion}
Let $\Pi$ satisfy \cref{V:orth} and let its range contain every deviatoric{-}symmetric tensor.
Then $\Pi$ satisfies \cref{V:korn} on every domain, with $\CK=\sqrt2$. Every member of
$\mathcal{F}$ qualifies, with sharp constants
$\CK(\mathbb{I})=1$, $\CK(\SPi)=\CK(\DSPi)=\sqrt2$ and $\CK(\DPi)=\sqrt{d/(d-1)}$.
\end{proposition}

\noindent The proof --- a pointwise Pythagoras identity chained with the $H^1_0$ identity
$\norm{\DSPi\nabla\bv}^2=\tfrac12\norm{\nabla\bv}^2+(\tfrac12-\tfrac1d)\norm{\nabla\cdot\bv}^2$
--- is given in~\cite{paper} and is not repeated here. What follows is the converse question:
\emph{how much} of \cref{V:korn} does that criterion capture?

\section{The sharp characterization}\label{sec:sharp}

The answer is that Korn compatibility is a purely algebraic property of $\Pi$, independent of
$\Omega$, and detectable on rank{-}one tensors alone.

\begin{definition}[Ellipticity constant]\label{def:cPi}
For a constant{-}coefficient linear map $\Pi$ on second{-}order tensors set
\begin{equation}\label{eq:cPi}
  \cP\ \coloneqq\ \min\bigl\{\,\bigl|\Pi(\bv\otimes\bk)\bigr|\ :\ |\bv|=|\bk|=1\,\bigr\},
\end{equation}
the minimum being attained, the set being compact and the map continuous.
\end{definition}

\begin{theorem}[Korn compatibility is equivalent to rank{-}one ellipticity]\label{thm:sharp}
Let $\Pi$ satisfy \cref{V:orth}. The following are equivalent:
\begin{enumerate}[label=\textup{(\roman*)},leftmargin=2.6em]
  \item\label{it:korn} $\Pi$ satisfies \cref{eq:PiKorn} on some bounded domain
  $\Omega\subset\mathbb{R}^d$;
  \item\label{it:allomega} $\Pi$ satisfies \cref{eq:PiKorn} on \emph{every} domain
  $\Omega\subset\mathbb{R}^d$;
  \item\label{it:cpi} $\cP>0$, i.e.\ the kernel of $\Pi$ contains no nonzero rank{-}one tensor.
\end{enumerate}
When they hold, the optimal constant is $\CK=\cP^{-1}$, and it is attained in the limit.
\end{theorem}

\begin{proof}
\ref{it:cpi}$\Rightarrow$\ref{it:allomega}. Let $\bv\in H^1_0(\Omega)^d$ and extend it by zero
to $\mathbb{R}^d$; the extension, still written $\bv$, lies in $H^1(\mathbb{R}^d)^d$ precisely
because the trace vanishes. Its distributional gradient transforms as
\begin{equation}\label{eq:fourier-grad}
  \widehat{\nabla\bv}(\bk)=\mathrm{i}\,\widehat{\bv}(\bk)\otimes\bk ,
\end{equation}
a rank{-}one tensor for each fixed $\bk$. Since $\Pi$ has real, constant coefficients it
commutes with the Fourier transform, and $\bk$ is real, so Plancherel's theorem gives
\begin{equation}\label{eq:plancherel}
\begin{aligned}
  \norm{\Pi\nabla\bv}^2
  &= \int_{\mathbb{R}^d}\bigl|\Pi\bigl(\widehat{\bv}(\bk)\otimes\bk\bigr)\bigr|^2\dd\bk \\
  &\ \ge\ \cP^2\int_{\mathbb{R}^d}\bigl|\widehat{\bv}(\bk)\bigr|^2\,|\bk|^2\dd\bk
  \ =\ \cP^2\,\norm{\nabla\bv}^2 ,
\end{aligned}
\end{equation}
the inequality being \cref{eq:cPi} applied to $\bk/|\bk|$ and to the real unit vectors
obtained by splitting $\widehat{\bv}=\boldsymbol{a}+\mathrm{i}\boldsymbol{b}$ --- legitimate
because $\Pi$ is real, so that
$\bigl|\Pi(\widehat{\bv}\otimes\bk)\bigr|^2
=\bigl|\Pi(\boldsymbol{a}\otimes\bk)\bigr|^2+\bigl|\Pi(\boldsymbol{b}\otimes\bk)\bigr|^2$
--- and then rescaled, the map $(\bv,\bk)\mapsto\Pi(\bv\otimes\bk)$ being bilinear. Hence \cref{eq:PiKorn} holds with
$\CK=\cP^{-1}$, on every domain and with no reference to its geometry.

\ref{it:allomega}$\Rightarrow$\ref{it:korn} is trivial.

\ref{it:korn}$\Rightarrow$\ref{it:cpi}. Suppose $\cP=0$, so that
$\Pi(\bv_0\otimes\bk_0)=0$ for some unit vectors $\bv_0,\bk_0$. Fix
$\varphi\in C_c^\infty(\Omega)$, $\varphi\not\equiv0$, and consider the oscillating profiles
\begin{equation}\label{eq:oscillating}
  \bv_\lambda(\bx)\ \coloneqq\ \lambda^{-1}\,\bv_0\,\varphi(\bx)\,
  \sin\bigl(\lambda\,\bk_0\!\cdot\!\bx\bigr)\ \in\ C_c^\infty(\Omega)^d ,
  \qquad \lambda>0 .
\end{equation}
Differentiating,
\begin{equation}\label{eq:oscillating-grad}
  \nabla\bv_\lambda
  = \underbrace{\bv_0\otimes\bk_0\;\varphi\,\cos(\lambda\,\bk_0\!\cdot\!\bx)}_{\text{order }1}
  \; + \; \underbrace{\lambda^{-1}\,\bv_0\otimes\nabla\varphi\,
    \sin(\lambda\,\bk_0\!\cdot\!\bx)}_{\text{order }\lambda^{-1}} .
\end{equation}
The first term is annihilated by $\Pi$ at every point, so
$\norm{\Pi\nabla\bv_\lambda}\le\norm{\nabla\varphi}\,\lambda^{-1}\to0$, whereas
$\norm{\nabla\bv_\lambda}^2\to\tfrac12\norm{\varphi}^2>0$ by the Riemann--Lebesgue lemma
applied to $\cos^2(\lambda\,\bk_0\!\cdot\!\bx)=\tfrac12(1+\cos(2\lambda\,\bk_0\!\cdot\!\bx))$.
Therefore $\norm{\nabla\bv_\lambda}/\norm{\Pi\nabla\bv_\lambda}\to\infty$, and since every
$\bv_\lambda$ already lies in $C_c^\infty(\Omega)^d\subset H^1_0(\Omega)^d$, no constant can
make \cref{eq:PiKorn} hold on $\Omega$.

Optimality of $\CK=\cP^{-1}$ follows by running \cref{eq:oscillating} at a minimizing pair
$(\bv_0,\bk_0)$ of \cref{eq:cPi}: there the leading term of \cref{eq:oscillating-grad} has
$|\Pi(\bv_0\otimes\bk_0)|=\cP$ exactly, so the quotient tends to $\cP^{-1}$.
\end{proof}

\begin{remark}[Why rank{-}one tensors are the whole story]\label{rem:whyrankone}
\Cref{eq:fourier-grad} is the reason: a gradient field has a rank{-}one Fourier symbol, so a
projector can only be blind to a velocity field if it is blind to a rank{-}one tensor. This
also explains the domain independence in \cref{thm:sharp}: the obstruction lives in the symbol,
not in $\Omega$. It is worth stressing where the boundary data enter: only \emph{once}, in the sufficiency
half, to extend $\bv$ by zero. The necessity half needs none --- its profiles are compactly
supported, so they defeat \cref{eq:PiKorn} on any space between $C_c^\infty(\Omega)^d$ and
$H^1(\Omega)^d$. On a domain with a Neumann portion the characterization therefore survives
in one direction only: $\cP>0$ remains necessary, whereas sufficiency becomes a genuine Korn
inequality, to which the Plancherel argument above does not apply.
\end{remark}

\section{The criterion is strictly weaker than the characterization}\label{sec:gap}

\begin{example}[An admissible projector outside $\mathcal{F}$]\label{ex:complement}
Let $\Pi=\mathbb{I}-\DSPi$, the projector onto the spherical{-}plus{-}skew part. Its range meets
the deviatoric{-}symmetric tensors only at the origin, so \cref{prop:criterion} says nothing
about it. Yet for unit $\bv,\bk$,
\begin{equation}\label{eq:complement-c}
  \bigl|(\mathbb{I}-\DSPi)(\bv\otimes\bk)\bigr|^2
  = 1 - \bigl|\DSPi(\bv\otimes\bk)\bigr|^2
  = \tfrac12\Bigl(1-(\bv\!\cdot\!\bk)^2\Bigr) + \tfrac1d(\bv\!\cdot\!\bk)^2 ,
\end{equation}
whose coefficient in $(\bv\!\cdot\!\bk)^2$ is $\tfrac1d-\tfrac12\le0$, so the minimum
sits at $(\bv\!\cdot\!\bk)^2=1$ and $\cP=d^{-1/2}>0$ (for $d=2$ the form is constant,
equal to $\tfrac12$, and every pair attains the minimum). By
\cref{thm:sharp}, $\mathbb{I}-\DSPi$ is admissible, with $\CK=\sqrt d$. Note that it is
admissible \emph{for the same reason} the criterion fails to see it: what matters is the
kernel, not the range.
\end{example}

\begin{example}[The two elementary projectors that are not admissible]\label{ex:failures}
The skew{-}symmetric part $\mathbf{T}\mapsto\tfrac12(\mathbf{T}-\mathbf{T}^{\mathsf T})$ has
$\cP=0$: its kernel contains every symmetric rank{-}one tensor $\bw\otimes\bw$, and the
degeneracy is realized inside $H^1_0(\Omega)^d$ by the gradient fields $\bv=\nabla\varphi$,
whose gradient is the (symmetric) Hessian of $\varphi$. The spherical part
$\mathbf{T}\mapsto\tfrac1d(\operatorname{tr}\mathbf{T})\mathbf{I}$ also has $\cP=0$: its kernel
contains every traceless rank{-}one tensor, such as $\be_1\otimes\be_2$, and the degeneracy is
realized by the divergence{-}free fields. Neither, therefore, can serve as a viscous projector
--- as one would expect on physical grounds, since neither controls the full velocity gradient.
\end{example}

\begin{example}[A one{-}parameter family interpolating between the two sides]\label{ex:tilted}
Admissibility is not a discrete alternative: there are admissible projectors with $\cP$
arbitrarily close to $0$. Take $d=2$ and the orthonormal tensor basis
\begin{equation}\label{eq:basis2d}
  \tfrac{1}{\sqrt2}\mathbf{I},\qquad
  \mathbf{E}_1=\tfrac{1}{\sqrt2}\begin{pmatrix}1&0\\0&-1\end{pmatrix},\qquad
  \mathbf{E}_2=\tfrac{1}{\sqrt2}\begin{pmatrix}0&1\\1&0\end{pmatrix},\qquad
  \mathbf{W}=\tfrac{1}{\sqrt2}\begin{pmatrix}0&1\\-1&0\end{pmatrix},
\end{equation}
and let $\Pi_\theta$ be the orthogonal projection onto the plane spanned by
$\cos\theta\,\tfrac{1}{\sqrt2}\mathbf{I}+\sin\theta\,\mathbf{E}_1$ and
$\cos\theta\,\mathbf{W}-\sin\theta\,\mathbf{E}_2$. A direct computation gives, for unit
$\bv$ and $\bk=(\cos\beta,\sin\beta)$,
\begin{equation}\label{eq:tilted-quad}
  \bigl|\Pi_\theta(\bv\otimes\bk)\bigr|^2
  = \tfrac12\bigl(1+\sin 2\theta\,\cos 2\beta\bigr),
\end{equation}
independent of $\bv$, whence
\begin{equation}\label{eq:tilted-c}
  \cP^2(\Pi_\theta)=\tfrac12\bigl(1-|\sin 2\theta|\bigr) .
\end{equation}
So $\Pi_\theta$ is admissible for every $|\theta|<\pi/4$, with $\CK\to\infty$ as
$|\theta|\to\pi/4$. At $\theta=0$ one recovers $\Pi_0=\mathbb{I}-\DSPi$ of
\cref{ex:complement}, consistently with $\cP^2=\tfrac12=d^{-1}$ at $d=2$.
\end{example}

\section{Consequence for the stabilization parameters}\label{sec:tau}

The article designs its stabilization parameters from the Fourier symbol of the differential
operator. The viscous block is $(2\alpha\nu/h^2)\,\mathbf{T}_{\Pi}(\bk_0)$ with
$\mathbf{T}_{\Pi}(\bk)_{ab}\coloneqq P_{aibj}k_ik_j$, whose quadratic form is
\begin{equation}\label{eq:symbol-quad}
  \bv^{\mathsf T}\mathbf{T}_{\Pi}(\bk)\,\bv=\bigl|\Pi(\bv\otimes\bk)\bigr|^2 ,
\end{equation}
by idempotency and self{-}adjointness. Two bounds follow immediately, and they are of quite
different quality:
\begin{itemize}[leftmargin=1.6em]
  \item \emph{Upper bound, uniform.} Non{-}expansiveness \cref{V:orth} gives
  $\mathbf{T}_{\Pi}(\bk_0)\le|\bk_0|^2\,\mathbb{I}_d$ for every admissible $\Pi$.
  \item \emph{Lower bound, not uniform.} \Cref{eq:cPi} and \cref{eq:symbol-quad} give
  $\mathbf{T}_{\Pi}(\bk_0)\ge\cP^2\,|\bk_0|^2\,\mathbb{I}_d$, and $\cP$ is positive but,
  by \cref{ex:tilted}, not bounded below over the admissible class.
\end{itemize}
For $\Pi\in\mathcal{F}$ the optimality in \cref{thm:sharp} turns $\CK\le\sqrt2$ into
$\cP^2\ge\tfrac12$, upgrading the lower bound to
$\mathbf{T}_{\Pi}(\bk_0)\ge\tfrac12|\bk_0|^2\mathbb{I}_d$, so the $O(1)$ factor the article
discards when simplifying $\tau_{\nu,1}^{-1}$ lies in the uniform window $[1,2]$ --- it equals
$2-2/d$ for $\DSPi$ and exactly $2$ for $\mathbb{I}$, $\SPi$ and $\DPi$. For a general
admissible projector the window degrades to $[2\cP^2,2]$: still $O(1)$ for each fixed $\Pi$,
and still absorbed into the algorithmic constant $c_1$, but no longer uniform over the class.

\begin{remark}[The practical reading]\label{rem:practical}
Nothing here obstructs using an admissible projector outside $\mathcal{F}$: the a priori
analysis of~\cite{paper} is stated under \cref{V:orth,V:korn} alone and goes through verbatim,
with $\CK=\cP^{-1}$ entering only the definiteness of the working norm. What is lost is
\emph{uniformity}: the algorithmic constant $c_1$ calibrated for $\mathcal{F}$ need not be
appropriate for a projector with small $\cP$, and the natural safeguard is to keep
$\cP$ bounded away from zero when a new viscous operator is proposed. By
\cref{eq:symbol-quad} that is a Legendre--Hadamard test,
\begin{equation}\label{eq:cPi-recipe}
  \cP^2=\min_{|\bk|=1}\ \lambda_{\min}\bigl(\mathbf{T}_{\Pi}(\bk)\bigr),
\end{equation}
a $d\times d$ eigenvalue problem minimized over the unit sphere --- \emph{not} a one-shot
eigenvalue computation on the $d^2\times d^2$ matrix of $\Pi$, whose spectrum is contained
in $\{0,1\}$ for every orthogonal projection and so carries no information about $\cP$.
\end{remark}

\section*{Machine verification}
All the algebra above is checked symbolically, for $d=2$ and $d=3$, by the script
\texttt{projector\_\allowbreak algebra\_\allowbreak verification.py} under \texttt{proof\_\allowbreak verification/\allowbreak sympy/}:
section~C computes
$\cP$ for $\mathbb{I}$, $\SPi$, $\DPi$, $\DSPi$, $\mathbb{I}-\DSPi$, the skew and the spherical
parts and confirms \cref{eq:complement-c} and the two vanishing constants of
\cref{ex:failures}; section~E reproduces \cref{eq:tilted-quad,eq:tilted-c} and exhibits
$\theta=-0.3$ and $\theta=-0.75$ as witnesses that the window $[1,2]$ genuinely fails outside
$\mathcal{F}$. The abstract chaining behind \cref{prop:criterion} is machine{-}checked in Coq, in
\texttt{ViscousProjector.v} under \texttt{proof\_\allowbreak verification/\allowbreak coq-formal/}.

\begin{thebibliography}{1}
\bibitem{paper}
G.~Casas, J.~Gonz\'alez-Us\'ua, R.~Codina and I.~de-Pouplana,
\emph{A stabilized finite element method for incompressible, inertial flows in inhomogeneous
porous media}, manuscript.
\end{thebibliography}

\end{document}
